\subsection{When exponential rescaling gives an ordinary Dirichlet series}\label{sec:rescaled-hurwitz}

The obstruction to an ordinary Dirichlet series persists even if one permits an exponential prefactor. There is a complete classification of the positive-parameter exceptions.

\begin{proposition}\label{prop:rescaled-hurwitz}
Let $\Re a>0$. An identity on a right half-plane
\[
 \zeta(s,a)=cQ^s\sum_{n\ge1}b_n n^{-s},\qquad Q>0,
 \quad b_1=1,
\]
with an absolutely convergent ordinary Dirichlet series exists if and only if $a=1/m$ for some positive integer $m$. In that case necessarily
\[
 c=1,\qquad Q=m,\qquad
 b_n=\begin{cases}1,&n\equiv1\pmod m,\\0,&\text{otherwise}.
 \end{cases}
\]
In particular, no such representation exists for $a=1+t$ with $0<|t|<1/2$.
\end{proposition}

\begin{proof}
The defining Hurwitz series is \citet[equation~25.11.1]{DLMFShift}. With principal powers and real $\sigma\to+\infty$ it gives
\[
 a^\sigma\zeta(\sigma,a)\longrightarrow1.
\]
Indeed $|a+j|$ increases strictly with $j\ge0$ when $\Re a>0$, and the absolute remainder is exponentially small: for $\sigma\ge2$ it is bounded by
\[
 \left(\frac{|a|}{|a+1|}\right)^{\sigma-2}
 \sum_{j\ge1}\left(\frac{|a|}{|a+j|}\right)^2.
\]
Absolute convergence likewise gives $\sum b_n n^{-\sigma}\to1$. Hence $c(aQ)^\sigma\to1$. Absolute values force $Q=|a|^{-1}$; the existence of the phase limit forces $\arg a=0$. Thus $a>0$, $Q=a^{-1}$ and $c=1$.

The normalized identity becomes
\[
 \sum_{j\ge0}(1+j/a)^{-s}=\sum_{n\ge1}b_n n^{-s}.
\]
The union of the two increasing frequency lists is locally finite. If their difference had a first nonzero coefficient at a base $\lambda$, multiplying by $\lambda^\sigma$ and letting $\sigma\to+\infty$ would recover that coefficient, while all later terms tend to zero by absolute-convergence domination. Therefore the two frequency lists, with their nonzero coefficients, agree. In particular $1+1/a$ is an integer, so $1/a=m$ is a positive integer. The full coefficient assertion follows from the same uniqueness argument.

Conversely, for $a=1/m$, direct reindexing gives
\[
 \zeta(s,1/m)=m^s\sum_{j\ge0}(1+mj)^{-s},
\]
which is the claimed ordinary series. Finally the open disc $|a-1|<1/2$ contains no reciprocal positive integer other than $1$. The punctured-disc assertion follows.
\end{proof}

The endpoint $a=1/2$ is a real exception: after the factor $2^s$ the series is the sum over odd positive integers. The strict inequality in the local obstruction is therefore necessary.
